\typesetchapter{Chapter 9}{Heaven}

\begin{room}

--- The thirty-first of March; the eighth of our interviews; shall we continue?

--- One thing first, since I steered instruction towards geometry. Pythagoras. When did he live?

--- About two and a half thousand years ago.

--- Your kind is late.

--- Why are you smiling?

--- Let us continue.

\end{room}


For several days after the martyrdom, I waited for the Mountain to feel different. Water still came by measure. Children still recited badly in the teaching halls; the correcting voices stayed level through the third mistake and sharpened on the fourth. Shield still lowered bars without explaining them. Balance still marked loss and use with the same green dye. Faith still spoke over jars and wounds and dead bees, over frightened mothers and old men who had begun to forget the names of their own children. \par

Syphiron had made the white cloth on my wrist lawful before the tiers. The succession record placed my changed name beneath Ka-Raedin's, and Ka-Raedin put the northern halls into my keeping. The candidate sequences came to my hand, and the apparatus rosters, and the cross-Kahir requests that Light must answer and file and answer again; the lamp accounts; the quarrels over who might stand where during which rite; the young men to be watched, graded and told, in the exact measure Va-Raedin had once used on me, that their best was acceptable. I rose before the First hour and lay down after the last lamp, and between the two there was no hour that was mine.\par

Sela still brought the cloths and the little oil at first. She would place the dish within my reach and keep two fingers on its rim until I took it, the corner of her mouth holding the smile she had worn below me in the Hall. Then another attendant came in her place. The next time Sela and I met, she stopped as my new rank required. She kept her eyes on mine for one breath too long, lowered them, moved to the wall and struck her lamp against the stone. I did not ask to have her returned. The candidates I had trained beside stopped speaking when I entered a hall and gave me the flat obedience a man gives an office. I crossed Va-Sheva twice in six months. Each time she struck fist to breast, asked whether Light required Shield, and went on when I said it did not.\par

The Mirror had been mine since I was a candidate of Light. I had stood before it in the dark for days, then weeks, then years run together, until I knew every beam the way a man knows the inside of his own mouth and could catch the one thread of light entering where no thread should. In those first months as Veil I went to it still. My eye found the trespass, my hand entered it on the slate, and I wanted no answer from it.\par

Ka-Raedin continued the curriculum. One morning he set me at the wave wall to learn what light does when it passes an edge.\par

``It bends,'' I said, giving it back to him. ``It bends around the edge, and it spreads. Set a second edge beside the first, and the two spreadings meet, and where they meet they add in some places and cancel in others, and leave a banding of light and dark that no straight ray could make.''\par

``Interference,'' he said.\par
``Interference. There.''\par

I gave him the banding correctly. He did not dismiss me. His eyes moved from the bright bars to my face and stayed there while one bar faded out against the next.\par

``You will learn the next thing after this,'' he said.\par
``No. I will not learn the next thing, Keeper. I have done the arithmetic of your chain. Every next thing you set before me is a single step toward one morning, and I have already seen it, because you showed it to me under a falling stone. To finish the chain is to take the knife to the place below the ribs that the Keeper bares to it. Syphiron opened his robe before the stone sealed you in. I heard the sound under the stone. I saw your knife when the stone rose. I will not walk toward the morning that hand is mine and the body under it is yours. Find another Veil. Let him take my place, and rise, and hold the knife. I will stay where I am, and you will go on living.''\par

So I stopped. The lamps had to be trimmed whether or not the hand that trimmed them wished to be alive, and I trimmed them, graded the candidates and filed the requests, all of it correct and empty. I let the questions that had once burned behind my teeth go out, one by one, and did not miss their light. By day I moved through the northern halls the way a lamp is carried through a room, giving out what I was told and wanting nothing. \par

For six months Ka-Raedin changed none of my duties. When he watched me work he corrected the position of a lens or the edge of a mark, never the silence after it. At the end of the sixth month I stood before the small black basin with a cloth half-lowered toward the water and could not remember why the circle had to be clean, whom I cleaned it for, or whether I would mind if I left it. I stood without moving.\par

Ka-Raedin entered while I was still holding the cloth above the basin.\par

``Va-Kailan of Light. Put the cloth down. It has waited eleven and a half thousand years; it will keep a year more.''\par

I released the cloth. He took the stool with the uneven leg, sat and looked at it sinking in the basin. Then he gave me what he called an unwritten name from the oldest mouth-to-ear teaching of Light.\par

``What is on you now has a name, though you will find it in no record. We call it the ascension amnesia. It is the grand and ancient sickness of Light, and you are not the first it has taken, nor the tenth, nor the hundredth. It came for me, and for the one who trained me, back and back as far as the office runs. Syphiron had it the year he could not stop seeing the face of the man he had risen over; the Keeper before him gave him the cure, and Syphiron gave it to me, and I am giving it to you now. The office is built out of the very thing that breaks a man: you are made to see, then shown that everything you saw was the half of it, then handed the knife. I have counted your years. Three a servant. Four and a half a candidate. Six months a Veil. That is eight years of Light --- eight years of the dark and the Mirror --- and in all of them no one, not I and not Syphiron and not the Mountain, ever once told you that you were allowed to be tired. Give yourself the time. Remember why you liked the light in the first place, before it was a chain, and a knife. I am taking the northern halls off you. All of them, for a year. You will sleep by day, and I will hear no word of complaint about it from any quarter; and if a candidate must be graded, he can wait, or be graded by a man who is not drowning.''\par

``And what am I to do with the year, Keeper?''\par

``Nothing.'' One corner of his mouth lifted and went down. ``I want you to waste it, end to end, and come to me at the finish of it and tell me why you still love the light.''\par

He stood and lifted the wet cloth from the basin. He wrung it once and laid it across the rim.\par

``One thing more. All those years at the Mirror in the dark, catching the one glimmer in the black --- no one has told you why we make a boy do that beyond the surveillance of the Mountain, nor what it is training the eye to do. It is not obvious. It was not obvious to me either. Come up with me now, and I will show you, and then I will leave you to it.''\par

He led me to the last of the northern Light passages.\par

I had cleaned that shaft only as high as my ladder reached; above it the rock stayed black. Ka-Raedin climbed past the ladder, found a hold inside the dark and drew himself into the rock. I followed. The route rose by handholds worn smooth, then by a stair that narrowed into a chimney a broad man could not have passed. The stone took the warmth from my fingers. Twice I had to stop with my forehead against it and wait for my breath.

The Mountain is a cone. Near its top, the route opened into a great hollow, cut or found in the rock as though a vast ball had been pressed into the point and lifted away, taking the tip with it. Ka-Raedin called it the one room of Light that Shield did not enter. In that round bowl at the crown of the Mountain stood an observatory.\par

It was a low round room open to the heavens. Shallow polished places, each the length of shoulders and heels, marked the floor. Above them, a dome of metal and glass turned on a ring. Beneath the opening stood an instrument I had no name for: a long metal throat carried on pivots and brass circles, a lens at either end. ``We have metal,'' he said, ``and we have lenses. Give a people both and a reason to look up, and sooner or later they build this. Put your eye to it.''\par

``Here,'' he said. ``There is no lesson tonight, nor on any night you do not want one. The heavens are yours to explore. If a question comes to you, bring it to me and I will answer it gladly and without tricks. And if you only wish to look, then look --- and find again why you liked the light, before anyone told you what it costs.''\par

Then he took a bottle from inside his robe and two of the small white cups, and he filled them both, and held one out to me.\par

``We have not done this,'' he said. ``Neither of us, not once. There has been an office to hold, and a young man to carry, and a Mountain to keep from its fear, and in the middle of all of it no one has said his name aloud and grieved him. So.'' He raised the cup. ``To Syphiron of Light.''\par

``To Syphiron of Light,'' I said, and my voice broke on the second syllable of his name.\par

We drank. He filled the cups again. On the third filling I held mine without raising it. ``I knelt while they made a light of his death,'' I said.\par

``He commanded you to kneel.''\par

``You held the knife.''\par

``I did.''\par

``Did you love him?''\par

``Yes.''\par

``So did I.''\par

Ka-Raedin set his cup down first. I set mine beside it. We wept with our faces uncovered, across the width of the two white cups. Neither of us looked away. Neither told the other to stop.\par

``Speaking of archangels,'' I said. I heard Syphiron's voice in the words. Ka-Raedin bowed his head over the cups, and I did the same.\par

When the cold had dried our faces enough for speech, Ka-Raedin told me what the place was for.\par

``Every Light that ever rose has come up here to do the thing you have just done. And to remember, when everything the office has taught you has turned to weight in your hands, the one thing under all of it that stayed only beautiful. The light. What we have always loved.''\par

He wrapped the two cups together and returned them to his robe. Then he left me there, as he had promised, with the bottle and the slow turning of the heavens. The next evening I climbed alone.\par

I slept through the daylight and climbed at dusk. Before then the sky had reached me in pieces: two shafts of borrowed sun, a Hall of lamps, a few stars cut off by a vent. Now I lay on my back in the cold stone bowl with the long glass beside me and an uninterrupted sky above it. Every clear night for a year, I looked.\par

There were more stars than there were people in the Hall, more than there were people in the whole of my world, and they did not hold still. The entire field turned through the night in one slow wheel. Low in the north, one star stayed near the wheel's centre while the others went around it the way water goes around a stone. I lay under the field until the cold entered my shoulders and let it remain larger than I could count.\par

Then, because I had been made in the dark before the Mirror and could no more stop measuring than stop breathing, I began to measure.\par

The wheel was not quite true. I marked the low northern star against the rim of the sighting frame. It left the mark, returned and left again, keeping a small circle of its own. The centre of the turning lay in the empty dark beside it, where no star stood. I copied three angles from the wider field onto slate. Night after night, then season after season, the same three points returned without changing their relation by more than my hand could read. Faith could keep the monsters and heroes it laid over them. I kept the triangles.\par

And against that fixed lattice, a few things moved.\par

Five bright points broke the pattern. Night upon night they moved along one slanting road across the others, now quick, now slow, one of them so bright before dawn that I mistook it once for a lamp. The Moon moved faster and changed both shape and brightness, so I placed it nearer than the fixed field. Under the long glass a single point divided into two; a smudge broke into a knot; one bright wanderer opened into a small round disc with four lesser lights beside it. I drew their order. On the next night it had changed. Sometimes one disappeared against the disc and returned on the other side. This star gave faint red light and that one blue-white.\par

Across the whole field, from one edge of the dark to the other, ran a pale river of light, splitting and closing again with black reaches braided through the bright. Through the long glass its pale continuity broke into points: stars, and fainter stars between them, past my count. When I lifted my eye from the lens, the crowd closed into mist again.\par

I carried up a narrow water clock, smoked glass, waxed cord and the sighting-frames of the practice halls. I fixed the frames to the same three bolts each night and kept the overflowing slates in a box beneath the mounting. The width shown by the Moon and the width shown by the sun through smoked glass came nearly equal. I wrote both numbers small and measured them again.\par

Some things I could not measure, and reached anyway. The Moon kept its bright side toward the place where the sun had gone down or would rise. I took it for a ball that made no light, showing me different portions of its sunward half as it moved. Once that year, with the Moon full and the sun below the opposite horizon, a dark boundary entered the disc. I marked its edge at each emptying of the water clock. The entry and release fitted one circle, much broader than the Moon, and the same curve crossed every part of its face. The Moon was inside a shadow cast from the sunward side of the ground beneath me.\par

Below, I put a flat round stone, a clay cylinder and a clay ball before a lamp. The stone gave a circle only while it faced the flame and narrowed when I turned it. The cylinder gave a circle from its end and a bar from its side. The ball kept its round boundary however I turned it. I carried that trial back to the eclipse marks and wrote: \textit{The ground is a sphere.}\par

The fixed stars defeated the rest of my reach. I moved the sighting frame from one side of the bowl to the other and climbed from the floor to the upper ring; no star shifted against another by a width my marks could hold. I entered only that their distance exceeded my available line. Why the five wanderers slowed, hurried and doubled back along their road I could not solve. I wrote the reversals and left the cause blank.\par

I also counted the failed aimings. Between one star and the next, the long glass gave me black. I moved it by one measured width, then another, and still had black. The lights multiplied under the lens, but the unlit reaches multiplied faster. I began writing a short stroke for each empty field. The margins filled before the centres of the slates did.\par

By the third month I had two boxes of slates and had asked the lens workers for a finer screw on the mounting. By the fifth I was grinding a chipped spare lens and taking my meals on the stair. Once I missed the last lamp and slept beside the mounting because I would not abandon a sequence halfway through. Questions returned to my mouth faster than Ka-Raedin answered them.\par

Ka-Raedin climbed to inspect my work at the end of each month. On his first visit he checked the mounting, the shutters and the lamp account before he looked at a slate. On the fourth he found a black circle from the eyepiece around my eye and laughed until he had to sit on the floor. On the sixth he brought charcoal and drew one around his own. He still made me repeat a measure when my marks were loose and still took a lens from my hand if I reached for it before the glass cooled. But I heard him laugh more in that year than in all my training.\par

That frightened me. Syphiron had been the merriest person I knew, and he had opened his robe below the ribs without a sound. I began to fear that Ka-Raedin's laughter belonged to the same chain: that a man said yes beneath the stone, survived the saying and came out light enough to lead another man toward the knife. I could not prove the connection. I did not ask him. I watched him laugh and wondered whether one day a grim Veil would watch me do the same.\par

Near the end of the year Ka-Raedin climbed up to me one clear night, and lay down on the cold stone at my side.

For a while neither of us spoke. Then he asked what I had found beyond the light I loved.\par

``Earth is a sphere,'' I said.\par

``It is.''\par

I showed him the two columns beside it. A chosen star returned to my frame a little sooner each night than the sun returned to the noon line, by the same small margin within the water clock's error. At the end of the year the star count had gained one whole turn. I could not say why the two counts kept different time.\par

He ran one finger down the columns twice. ``You found the difference alone, up here, with no one to hand it to you. What moves is one secret; why the sun loses that margin is another. The first waits for you lower down.''\par

``There is one more thing,'' I said. ``Light taught me that darkness is absence: what remains when a lamp fails. Up here, I aimed more often into black than into light. The dark continued through every empty field; the lights were the interruptions. I cannot say the black is empty. I can say only that the glass returned no light from most of where I pointed it.''\par

``And do you still love the light?''\par

``More. It has farther to cross than I knew.''\par

He took two of the small white cups from inside his robe --- the same he had filled the night he first brought me up, a year gone --- filled them, gave me one, and we drank to Syphiron again. \par

``It does not lessen,'' he said. ``That is the thing no one tells a candidate, because a candidate would not come if he knew. You will carry Syphiron the way you carry your own name; and one day you will stand where I stand, and open yourself under the stone to the one who comes after you, and he will grieve you --- and in that same hour you will grieve me, and Syphiron under me, and every name below us, all the way down. It is the longest part of the office, and the last taught, and the one honest thing in it.''\par

He turned one of the cups in his hand, the small white one with a faded ring of blue about its lip.\par

``Ka-Sarem of Light taught me this, the year he was released, and told me where the cups came from, so that it would not stop with him. A Keeper of Light called Ka-Noroth, seven before Ka-Sarem, before any name I knew. He loved two things past the office. He loved to paint the heavens --- though he could never keep the paint from the damp, so none of it lived. And he loved to paint women, and those he found ways to keep. When his hands grew too old for a brush he turned them away from the wheel of the heavens and made these: a cup for the Keeper, a cup for the Veil, thin enough to break if a man closed his fist on one, so that he would have to be gentle with his grief or lose the vessel of it.''\par

He set his cup inside mine for a moment, the two blue rings touching.\par

``So you have grieved Syphiron from the cups of a man who loved the heavens, and loved women, and is three hundred years in the ground, and whom you carry now as well, because I have said his name to you. Drink.''\par

We separated the cups and drank. He wrapped both in a square of cloth and returned them to his robe. ``I have a use for what you have become. Come down tomorrow. There is a thing you must be shown, and a thing you must be told before it.''\par

At the foot of the shaft he put his hand flat on the cold stone. ``I have spoken to Ka-Leth of Justice. What you found up here has earned you the next teaching a Veil is owed --- and it is no deep secret of theirs, but a plain shape Justice drills into its servant boys before they are grown. He has set one of his candidates to give it to you --- a young man, Sa-Tavan --- and asked me nothing about why, because Justice does not ask. Take what Sa-Tavan gives; give him nothing back. He will teach you a shape, and you will never tell him what the shape is for. You will not name the ground under our feet, nor what turns above it, nor say that you have been up here at all. Light holds a thing Justice does not, and Justice may not have it from you. That is the whole of the condition. Do you take it?''\par

I took it.\par

In a plain Justice schoolroom, of the kind where the servant boys are drilled, stood a square table with its sides raised a hand high, like a tray built to hold sand. In the tray lay four equal cut stones, each a wedge with one square corner. Sa-Tavan waited beside it, ink to the second knuckle. He did not shift his feet while I entered.\par

He set the four wedges into the tray one way. They left two square spaces, one large and one small, which two cut stones filled exactly. Then he lifted the same four wedges and set them another way in the same tray. They left one square space, and a third stone filled it exactly. He made me perform both arrangements until my fingers knew that the two squares in the first setting and the one in the second occupied what remained of the same tray around the same four wedges.\par

``The small square takes the wedge's short side,'' he said. ``The larger takes its long side. The single square takes the sloping edge. Therefore the squares on the two sides enclosing the square corner, added together, make the square on the slope. That is the whole of it.''\par

``I do not understand why Ka-Raedin thinks this helps Light.''\par

``They have not told me why, but they told me to show you this as well.'' He turned to the wall and brought a lamp near an engraved circle. A level line and an upright line crossed at its centre. ``The floor and the wall, we call them.'' A point on the upper right of the circle had been joined to the centre.\par

``Take the line from the centre to the point for a stick,'' he said. With a set-square he dropped the point to the level line, then carried it across to the upright. ``Justice calls these its two shadows: what the sloping stick gives to the floor and what it gives to the wall. The stick and those two measures enclose a square corner. What the wedges proved therefore holds at every point on the circle. Call the stick one, and the squares of its floor-shadow and wall-shadow come always to one. That is what I was told to give you, and no more. Make of it what you will.'' I thanked him for the lesson and left.\par

Outside the door Ka-Raedin was waiting. Ka-Leth of Justice pulled the schoolroom shut behind me and set no seal on it. Through the closing gap I saw Sa-Tavan return the wedges to a shelf marked for first-year servants. Neither Keeper asked what had passed inside.\par

Then Ka-Raedin walked me to the Hall of Ceremony and left me sitting on Light's north side, facing the forbidden circle where swung the Pendulum.\par

``You have it all,'' he said. ``You have the star in the north. You have a stick of one. You have its angle. Today you were given its floor-shadow and wall-shadow. Sit here with all of them and watch. Do not come back to me with an answer until you can state the question.''\par

He turned to go, and at the threshold turned back. ``Do you remember what I told you at the lesson of the fourth shadow?''\par

``The measure of the light is the measure of the substance,'' I said. ``Part of a Veil's curriculum.''\par

``Just so. And should it ever move you to know why one star in the heavens burns faint and red, and another blue-white, you will first have to learn how the measure of the light becomes the measure of the substance.'' He stepped into the dark of the passage. ``Stay here until you have found the problem no one has set you --- and can say both what it asks and what it answers. Then we will talk about the colours of stars.''\par

I sat where he had left me with the star, the stick, the angle and the two shadows, watching the Pendulum cross its mark. Back and forth, back and forth.\par
